% Zdroj: PDF priklady-podzim-2023.pdf, fyzická str. 1-2. Značka: společné okruhy. Zařazeno: I4.
\section{Architektura počítačů a operačních systémů}
\szzotazka{podzim 2023}{1}{Architektura PC a OS (race condition, kritická sekce)}

\noindent\textbf{Zadání.}\par
Mějme následující třídu v~pseudokódu:

\begin{lstlisting}[language=C++,numbers=left]
class LimitedStack {
public:
        LimitedStack(int md) {
            currentDepth = 0;
            maxDepth = md;
            stack = new int[maxDepth];
        }

        int getDepth() {
            return currentDepth;
        }

        bool pushValue(int v) {
            if (currentDepth >= maxDepth)
                return false;
            stack[currentDepth] = v;
            currentDepth++;
            return true;
        }

private:
        int      maxDepth;
        int      currentDepth;
        int[]    stack; // int *stack; pro C++
};
\end{lstlisting}

Dále předpokládejme, že máme multiprocesorový systém a~program si spustí několik vláken,
která budou mít všechna přístup ke sdílené proměnné typu \texttt{LimitedStack}. Tato vlákna
budou volat podle své potřeby funkce \texttt{getDepth} a~\texttt{pushValue}.

\begin{enumerate}
\item Je možné, že v~programu nastane jev zvaný \textit{race condition}? Pokud ano,
  napište rozsahy řádek nebo vyznačte přímo v~uvedeném kódu řádky, které představují
  kritickou sekci.
\item Je možné, aby nějakým způsobem nastala situace, kdy při volání funkce
  \texttt{pushValue} přeteče zásobník (program se bude pokoušet zapisovat mimo alokované
  pole)? Zdůvodněte.
\item Pokud se domníváte, že v~uvedeném kódu může nastat race condition, pokuste se tento
  problém odstranit úpravou programu (včetně možného přidání nějakých nových řádek).
\end{enumerate}
Kde je to důležité, ve své odpovědi uvažujte jazyk C\#, C++ nebo Java (a~vaši volbu vyznačte).

\medskip
\noindent\textbf{Náčrt řešení.}\par
\begin{enumerate}
\item Ano, 10 a~14-17.
\item Ano, T1 i~T2 vykonají řádek 14 souběžně. Přečtou shodnou proměnnou
  \texttt{currentDepth}, která má hodnotu $\texttt{maxDepth}-1$, takže test selže.
  Následně po testu je T2 plánovačem zastaveno, T1 funkci dokončí a~zvětší
  \texttt{currentDepth}. Poté je T2 znovu naplánováno a~při zápisu hodnoty použije
  zvětšený index sahající za pole.
\item Záleží na jazyku. V~Javě triviálně použitím \texttt{synchronized} metod. C\# buď
  použije \texttt{lock} na nějaký vedlejší objekt nebo se použije
  \texttt{[MethodImpl(MethodImplOptions.Synchronized)]}. C++ použije třeba \texttt{Mutex}
  a~buď explicitní \texttt{lock} a~\texttt{unlock} nebo lépe \texttt{lock\_guard}. Pozor
  na odemknutí před každým \texttt{return}.
\end{enumerate}

\medskip
\noindent\textit{Doplňující vysvětlení (nad rámec oficiálního náčrtu):} Řádek~10 patří do kritické
sekce proto, že čte \texttt{currentDepth}, kterou jiné vlákno právě mění; řádky 14--17 čtou i~zapisují
\texttt{currentDepth} a~zapisují do \texttt{stack}. Samotné \texttt{currentDepth++} navíc není atomické
(načti, přičti, ulož), takže se dvě souběžná zvýšení mohou slít v~jedno. Zámek musí obepnout
\emph{celý} úsek od testu na řádku~14 po inkrementaci: kdyby chránil jen zápis, dvě vlákna by dál
mohla projít testem se stejnou hodnotou. Možná úprava v~C++:
\begin{lstlisting}[language=C++]
std::mutex mtx;   // nova clenska promenna

int getDepth() {
    std::lock_guard<std::mutex> g(mtx);
    return currentDepth;
}
bool pushValue(int v) {
    std::lock_guard<std::mutex> g(mtx);   // odemkne se pri kazdem return
    if (currentDepth >= maxDepth)
        return false;
    stack[currentDepth] = v;
    currentDepth++;
    return true;
}
\end{lstlisting}
